\begin{table}[htbp]\centering
\begin{threeparttable}
\caption{Taille d'effet minimale détectable par permutation des cohortes}\label{tab:mde}
\small
\begin{tabular}{lrrr}
\toprule
Hypothèse & É.-t. des ATT placebo & MDE (80 \%, 5 \%), \% & Permutations \\
\midrule
H1 commune, log(naissances/1 000 f. 15-44) & 0.0030 & 0.84 & 200 \\
H2b département, log(naissances/1 000 f. 25-39), bascule D3 $\geq$ 50 \% & 0.0047 & 1.31 & 200 \\
H2 département, log(naissances/1 000 f. 15-19) & 0.0210 & 5.89 & 200 \\
H2 département, log(naissances/1 000 f. 20-24) & 0.0100 & 2.81 & 200 \\
H2 département, log(naissances/1 000 f. 25-29) & 0.0065 & 1.81 & 200 \\
H2 département, log(naissances/1 000 f. 30-34) & 0.0059 & 1.65 & 200 \\
H2 département, log(naissances/1 000 f. 35-39) & 0.0071 & 1.99 & 200 \\
\bottomrule
\end{tabular}
\begin{tablenotes}\footnotesize
\item Source : scripts/04\_sample\_mde.py. ATT statique TWFE sur le log du taux, cohortes permutées entre unités ; MDE = (1,96 + 0,84) $\times$ écart-type placebo.
\end{tablenotes}
\end{threeparttable}
\end{table}
